\chapter{Linear Dynamical Systems}
\label{ch:linsys}

\begin{objectives}
\begin{itemize}
  \item Write physical systems in continuous- and discrete-time state-space form.
  \item Discretize deterministic inputs and white process noise exactly and approximately, including Van Loan's method.
  \item Understand stability, and why inertial errors grow polynomially.
  \item Define observability, test it with the observability matrix, and compute unobservable directions.
  \item Understand what observability means for the covariance of a Kalman filter.
\end{itemize}
\end{objectives}

The Kalman filter is built on a linear state-space model. This chapter collects the linear systems theory it needs: how to write a model, how
to turn a continuous-time model into the discrete-time form a computer processes, and how to decide whether the measurements can determine the
state at all.

% ==============================================================================
\section{State-space models}

A continuous-time linear \term{state-space model} has the form
\begin{equation}
  \dot x(t) = F(t)\,x(t) + B(t)\,u(t) + G(t)\,w(t), \qquad z(t) = H(t)\,x(t) + n(t),
  \label{eq:ct-model}
\end{equation}
with state $x$, known input $u$, white process noise $w$ (density $Q_c$), measurement $z$ and measurement noise $n$. The state is the smallest set of
numbers whose present value, together with future inputs and noises, determines the future. Its discrete-time counterpart is
\begin{equation}
  x_{k+1} = \Phi_k\,x_k + \Gamma_k\,u_k + w_k, \qquad z_k = H_k\,x_k + n_k,
  \qquad w_k \sim \N(0, Q_k),\ n_k \sim \N(0, R_k).
  \label{eq:dt-model}
\end{equation}

\begin{example}[Models used in these notes]
\label{ex:state-models}
\begin{itemize}
  \item \emph{Constant velocity}: $x = (p, v)$, $\dot p = v$, $\dot v = w$ (white acceleration). $F = \begin{bmatrix}0 & 1\\0 & 0\end{bmatrix}$, $G = \begin{bmatrix}0\\1\end{bmatrix}$.
  \item \emph{Random-walk bias}: $\dot b = w_b$. $F = 0$, $G = 1$.
  \item \emph{Inertial error model} (one axis, preview of \cref{ch:errordyn}): $\delta\dot p = \delta v$, $\delta\dot v = g\,\delta\theta - \delta b_a - n_a$,
        $\delta\dot\theta = -\delta b_g - n_g$, plus random-walk biases. A chain of integrators driven by noise.
\end{itemize}
\end{example}

The discrete-time solution follows by repeated substitution:
\begin{equation}
  x_k = \Phi(k, 0)\,x_0 + \sum_{j=0}^{k-1}\Phi(k, j+1)\big(\Gamma_ju_j + w_j\big), \qquad \Phi(k, j) = \Phi_{k-1}\Phi_{k-2}\cdots\Phi_j,
\end{equation}
the discrete analog of the variation-of-constants formula (\cref{prop:stm}).

% ==============================================================================
\section{Discretization}
\label{sec:discretization}

Physical laws are continuous, computers are discrete. Over a step $\Dt$ in which $F$, $B$ and $G$ are constant, the solution of \eqref{eq:ct-model}
(\cref{thm:lti-solution}) is
\begin{equation}
  x(t + \Dt) = \underbrace{e^{F\Dt}}_{\Phi}\,x(t) + \int_0^\Dt e^{Fs}B\,u(t + \Dt - s)\,\dd s + \underbrace{\int_0^\Dt e^{Fs}G\,w(t + \Dt - s)\,\dd s}_{w_k},
\end{equation}
which has exactly the discrete form \eqref{eq:dt-model}.

\paragraph{Deterministic inputs.} If the input is held constant over the step (a \emph{zero-order hold}, which is what a sampled IMU reading
means), the input term is $\Gamma u_k$ with $\Gamma = \int_0^\Dt e^{Fs}\,\dd s\,B$. Both $\Phi$ and $\Gamma$ come from one matrix exponential,
\begin{equation}
  \exp\!\Big(\begin{bmatrix} F & B \\ 0 & 0\end{bmatrix}\Dt\Big) = \begin{bmatrix} \Phi & \Gamma \\ 0 & I\end{bmatrix},
\end{equation}
which can be verified by noting that both sides solve the same matrix ODE in $\Dt$.

\paragraph{Process noise.} The discrete noise $w_k$ is a weighted integral of white noise, so it is zero-mean Gaussian with covariance (by the
same calculation as in \cref{prop:lyapunov})
\begin{equation}
  Q_d = \int_0^\Dt e^{Fs}\,GQ_cG\T\,e^{F\T s}\,\dd s .
  \label{eq:qd-exact}
\end{equation}

\subsection{Approximations}

Expanding to first order in $\Dt$ gives the approximations used by most real-time filters and frozen in AGENTS.md F2:
\begin{equation}
  \Phi \approx I + F\Dt, \qquad Q_d \approx GQ_cG\T\,\Dt .
  \label{eq:disc-first-order}
\end{equation}
Both are accurate to $O(\Dt^2)$ per step. The units confirm the second: a noise density $\sigma_a$ in $\mathrm{m/s^2/\sqrt{Hz}}$ gives $\sigma_a^2\Dt$ in
$\mathrm{m^2/s^2}$, a velocity variance. A second-order option adds $\tfrac12F^2\Dt^2$ to $\Phi$ and $\tfrac12(FGQ_cG\T + GQ_cG\T F\T)\Dt^2$ to $Q_d$.

\begin{example}[Double integrator]
\label{ex:double-integrator-qd}
For position and velocity driven by white acceleration of density $\sigma$, $F = \begin{bmatrix}0&1\\0&0\end{bmatrix}$, $G = \begin{bmatrix}0\\1\end{bmatrix}$, $Q_c = \sigma^2$.
\Cref{ex:double-integrator} gives $e^{Fs}G = [s;\ 1]$, so \eqref{eq:qd-exact} evaluates exactly to
\begin{equation}
  Q_d = \sigma^2\begin{bmatrix} \Dt^3/3 & \Dt^2/2 \\ \Dt^2/2 & \Dt \end{bmatrix}.
\end{equation}
The first-order approximation $Q_d \approx \diag(0, \sigma^2\Dt)$ puts no noise on position within a single step. Yet after $n$ steps of exact $\Phi$, the
accumulated position variance is $\sum_{k=1}^{n} (k\Dt)^2\sigma^2\Dt \to \sigma^2 t^3/3$ as $\Dt\to0$ (a Riemann sum), which is the exact answer of
\cref{prop:integrated-noise}. The per-step error of the approximation does not accumulate into a wrong growth law; it only introduces an $O(\Dt)$
relative error, negligible at 200~Hz.
\end{example}

\subsection{Van Loan's method}

When higher accuracy is wanted, \eqref{eq:qd-exact} can be computed exactly with a single matrix exponential~\cite{vanloan1978}.

\begin{proposition}[Van Loan]
\label{prop:vanloan}
\index{Van Loan's method}
Let $C = \exp\!\Big(\begin{bmatrix} -F & GQ_cG\T \\ 0 & F\T \end{bmatrix}\Dt\Big) = \begin{bmatrix} C_{11} & C_{12} \\ 0 & C_{22}\end{bmatrix}$. Then
$\Phi = C_{22}\T$ and $Q_d = \Phi\, C_{12}$.
\end{proposition}
\begin{proof}
For a block upper-triangular generator, the diagonal blocks of the exponential are the exponentials of the diagonal blocks, so $C_{22} = e^{F\T\Dt}$.
The off-diagonal block, as a function of $\Dt$, satisfies $\dot C_{12} = -FC_{12} + GQ_cG\T C_{22}$ with $C_{12}(0) = 0$, whose solution (variation of constants)
is $C_{12} = \int_0^\Dt e^{-F(\Dt - s)}\,GQ_cG\T\,e^{F\T s}\,\dd s$. Multiplying by $\Phi = e^{F\Dt}$ on the left, and using
$e^{F\Dt}e^{-F(\Dt - s)} = e^{Fs}$, gives $\Phi C_{12} = \int_0^\Dt e^{Fs}\,GQ_cG\T\,e^{F\T s}\,\dd s$, which is \eqref{eq:qd-exact}.
\end{proof}

For the 15-state inertial error model this costs a $30\times30$ matrix exponential per step, too expensive at 200~Hz but an excellent test oracle for the
closed forms of \cref{ch:errordyn}.

% ==============================================================================
\section{Stability}

A linear time-invariant system is \emph{stable} if its free response decays. For $\dot x = Fx$ the modes are $e^{\lambda_it}$ (\cref{thm:lti-solution}), so the
system is asymptotically stable when all eigenvalues of $F$ have negative real parts. For $x_{k+1} = \Phi x_k$ the modes are $\lambda_i^k$, so all eigenvalues of
$\Phi$ must lie strictly inside the unit circle; the two conditions correspond through $\lambda_\Phi = e^{\lambda_F\Dt}$.

Integrators have eigenvalues at zero (continuous) or one (discrete): they are \emph{marginally} stable. A single integrator holds its state; a chain of
integrators (a nilpotent $F$) makes initial errors grow \emph{polynomially}, like $t$, $t^2$, $t^3$. The error dynamics of inertial navigation are of this
kind (\cref{ex:state-models}), which is why IMU errors grow polynomially rather than exponentially, and why they never stop growing without aiding.

% ==============================================================================
\section{Observability}
\label{sec:observability-linear}

\subsection{Definition and the observability matrix}

\begin{definition}[Observability]
\index{observability}
A system is \term{observable} if the initial state $x_0$ can be determined from the inputs and the noise-free outputs over a finite time.
\end{definition}

Because the inputs are known, their contribution to the outputs can be subtracted, so observability depends only on $(\Phi, H)$. For a discrete
time-invariant system without noise, $z_k = H\Phi^kx_0$. Stacking the first $n$ outputs,
\begin{equation}
  \begin{bmatrix} z_0 \\ z_1 \\ \vdots \\ z_{n-1}\end{bmatrix} = \underbrace{\begin{bmatrix} H \\ H\Phi \\ \vdots \\ H\Phi^{n-1}\end{bmatrix}}_{\mathcal O}x_0 .
\end{equation}

\begin{theorem}[Observability rank test]
\label{thm:obs-rank}
\index{observability matrix}
The system is observable if and only if the \term{observability matrix} $\mathcal O$ has rank $n$. Its nullspace is the \term{unobservable subspace}: the
initial states that produce identically zero output forever.
\end{theorem}
\begin{proof}
If $\mathcal O$ has full column rank, $x_0$ is the unique solution of the stacked equations. If $\mathcal Ox_0 = 0$ for some $x_0 \ne 0$, then $H\Phi^kx_0 = 0$ for
$k < n$, and by the Cayley--Hamilton theorem ($\Phi^n$ is a linear combination of $I, \Phi, \dots, \Phi^{n-1}$) also for all $k \ge n$; such an $x_0$ is
indistinguishable from zero. More outputs never add rank beyond the first $n$.
\end{proof}

The unobservable subspace is invariant under $\Phi$: an unobservable error stays unobservable as the system evolves. For time-varying systems, the
same reasoning gives $\mathcal O = [H_0;\ H_1\Phi(1, 0);\ H_2\Phi(2, 0);\ \dots]$, and observability can depend on the trajectory, a fact that matters
greatly for VIO (\cref{ch:observability}).

\subsection{Examples}

\begin{example}[Position versus velocity measurements]
For the discretized constant-velocity model, $\Phi = \begin{bmatrix}1 & \Dt\\0 & 1\end{bmatrix}$.
\begin{itemize}
  \item Measuring position, $H = [1\ \ 0]$: $\mathcal O = \begin{bmatrix}1 & 0\\1 & \Dt\end{bmatrix}$ has rank 2. Velocity is observable, through the change of position.
  \item Measuring velocity, $H = [0\ \ 1]$: $\mathcal O = \begin{bmatrix}0 & 1\\0 & 1\end{bmatrix}$ has rank 1, with nullspace $(1, 0)$. Position is unobservable: shifting the whole
        trajectory changes no measurement. This is dead reckoning in one line.
\end{itemize}
\end{example}

\begin{example}[A bias that looks like the state]
Let $x = (p, b)$ with $\dot p = 0$, $\dot b = 0$ and $z = p + b$. Then $H = [1\ \ 1]$, $\Phi = I$, and $\mathcal O = \begin{bmatrix}1 & 1\\1 & 1\end{bmatrix}$ has nullspace $(1, -1)$: a
change of position compensated by an opposite change of bias is invisible. A sensor bias can be estimated only if something else breaks this
symmetry. In inertial navigation that something is motion: an accelerometer bias becomes separable from tilt once the vehicle rotates
(\cref{ch:strapdown}).
\end{example}

\begin{example}[An IMU along a line]
Let $x = (p, v, b_a)$ with $\dot p = v$, $\dot v = a_m - b_a$ ($a_m$ known) and $\dot b_a = 0$, and measure $p$. In continuous time,
$F = \begin{bmatrix}0&1&0\\0&0&-1\\0&0&0\end{bmatrix}$ and $H = [1\ 0\ 0]$, and the continuous-time observability matrix $[H;\ HF;\ HF^2]$ is
$\begin{bmatrix}1&0&0\\0&1&0\\0&0&-1\end{bmatrix}$, of full rank. Position measurements reveal velocity (one differentiation) and the accelerometer bias (two
differentiations): the reason a pose sensor makes all IMU states observable in Milestone 6.
\end{example}

\subsection{Duality and controllability}

The dual notion, \emph{controllability}, asks whether the inputs (or noises) can drive the state anywhere; its test matrix is
$[\Gamma,\ \Phi\Gamma,\ \dots,\ \Phi^{n-1}\Gamma]$, and $(\Phi, \Gamma)$ is controllable exactly when $(\Phi\T, \Gamma\T)$ is observable. For a Kalman filter,
controllability by the process noise means that every state is continually excited by noise, which keeps the filter from becoming certain
of anything and ignoring new data.

\subsection{What observability means for a filter}

For a time-invariant system that is observable from its measurements and controllable from its process noise, the Kalman filter covariance
converges to a unique, bounded steady state, whatever the initial covariance (\cref{ch:kf}). Along unobservable directions there is no such
bound: the covariance there can only grow with process noise or stay at its prior value, because no measurement reduces it. A filter must
respect this. If its linearized model makes an unobservable direction appear observable, it will \emph{reduce} the covariance along that direction
without justification, and become overconfident. This is precisely what happens to the standard EKF in VIO (\cref{ch:observability}).

\begin{takeaways}
\begin{itemize}
  \item State-space models collect everything needed to predict the future; discrete models follow from continuous ones by integrating over
        a step.
  \item $\Phi = e^{F\Dt}$, $Q_d = \int e^{Fs}GQ_cG\T e^{F\T s}\dd s$; to first order $\Phi \approx I + F\Dt$ and $Q_d \approx GQ_cG\T\Dt$. Van Loan computes both exactly.
  \item Chains of integrators are marginally stable: their errors grow polynomially.
  \item Observability: $\mathcal O = [H;\ H\Phi;\ \dots]$ must have full rank; its nullspace contains the directions the outputs cannot see.
  \item A filter may reduce uncertainty only in observable directions.
\end{itemize}
\end{takeaways}

\begin{exercises}
  \item Write the mass-spring-damper $m\ddot y + c\dot y + ky = u$ in state-space form and find its eigenvalues. When is it stable?
  \item Compute $\Gamma$ for the double integrator with a zero-order-hold acceleration input, and interpret its entries.
  \item Verify \cref{prop:vanloan} numerically for the double integrator.
  \item Find the unobservable subspace of the constant-velocity model with measurements of $p_1 - p_2$ for two independent targets.
  \item Show that the continuous-time system $\dot x = Fx$, $z = Hx$ is observable if and only if $[H;\ HF;\ \dots;\ HF^{n-1}]$ has full rank.
  \item A gyroscope measures $\omega + b_g$, and an inclinometer measures the integrated angle. Is the gyro bias observable? What if the inclinometer
        is replaced by a second gyroscope?
\end{exercises}
% Solutions to the exercises of Chapter 11.
\begin{solutions}
  \item With $x = (y, \dot y)$,
        \[
          \dot x = \begin{bmatrix}0 & 1\\ -k/m & -c/m\end{bmatrix}x + \begin{bmatrix}0\\ 1/m\end{bmatrix}u .
        \]
        The characteristic polynomial is $\lambda^2 + (c/m)\lambda + k/m$, so $\lambda = \big(-c \pm \sqrt{c^2 - 4km}\big)/2m$, or
        $\lambda = -\zeta\omega_n \pm \omega_n\sqrt{\zeta^2 - 1}$ with $\omega_n = \sqrt{k/m}$ and $\zeta = c/(2\sqrt{km})$. For $m > 0$ both roots have negative real parts
        exactly when $c > 0$ and $k > 0$ (a quadratic with positive coefficients): the system is asymptotically stable. Underdamped ($\zeta < 1$) it
        oscillates with decaying amplitude; with $c = 0$ the eigenvalues $\pm i\omega_n$ lie on the imaginary axis and the system is only marginally stable;
        with $k < 0$ (a negative spring) one eigenvalue is positive and it is unstable. This is the accelerometer model of \cref{ch:imu}.

  \item With $F = \begin{bmatrix}0 & 1\\0 & 0\end{bmatrix}$ and $B = \begin{bmatrix}0\\1\end{bmatrix}$, $e^{Fs}B = \begin{bmatrix}s\\1\end{bmatrix}$ (\cref{ex:double-integrator}), so
        \[
          \Gamma = \int_0^\Dt e^{Fs}B\,\dd s = \begin{bmatrix}\Dt^2/2\\ \Dt\end{bmatrix}.
        \]
        An acceleration $a$ held constant over the step adds $\tfrac12a\Dt^2$ to the position and $a\Dt$ to the velocity: the familiar kinematics of constant
        acceleration, exact for a zero-order-hold input. The block exponential $\exp\big(\begin{bmatrix}F & B\\0 & 0\end{bmatrix}\Dt\big)$, whose generator is nilpotent
        ($3\times3$ with cube zero), gives the same result after two terms.

  \item With $\Dt = 0.1$ and $\sigma = 1$, \code{scipy.linalg.expm} of $\begin{bmatrix}-F & GQ_cG\T\\ 0 & F\T\end{bmatrix}\Dt$ gives
        \[
          \Phi = C_{22}\T = \begin{bmatrix}1 & 0.1\\0 & 1\end{bmatrix}, \qquad
          Q_d = \Phi C_{12} = \begin{bmatrix}3.333\times10^{-4} & 5\times10^{-3}\\ 5\times10^{-3} & 0.1\end{bmatrix},
        \]
        which equals the exact $\sigma^2\begin{bmatrix}\Dt^3/3 & \Dt^2/2\\ \Dt^2/2 & \Dt\end{bmatrix}$ of \cref{ex:double-integrator-qd} to $10^{-18}$.

  \item With $x = (p_1, v_1, p_2, v_2)$, $\Phi = \blkdiag\big(\begin{bmatrix}1 & \Dt\\0 & 1\end{bmatrix}, \begin{bmatrix}1 & \Dt\\0 & 1\end{bmatrix}\big)$ and $H = [1\ \ 0\ \ {-1}\ \ 0]$, the rows
        of the observability matrix are $H\Phi^k = [1\ \ k\Dt\ \ {-1}\ \ {-k\Dt}]$, which span a two-dimensional space: $\rank\mathcal O = 2$. The unobservable subspace
        is
        \[
          \nullsp(\mathcal O) = \mathrm{span}\{(1, 0, 1, 0),\ (0, 1, 0, 1)\}:
        \]
        a common translation and a common velocity of both targets. The relative measurement reveals the relative position and the relative
        velocity, but nothing about where the pair is or how it moves as a whole: the one-dimensional version of the unobservable global position of
        VIO.

  \item For $\dot x = Fx$, $z(t) = He^{Ft}x_0$, and the derivatives of the output at $t = 0$ are $z^{(k)}(0) = HF^kx_0$. If $\mathcal O_c = [H;\ HF;\ \dots;\ HF^{n-1}]$ has
        full column rank, the noise-free output over any interval determines these derivatives and hence $x_0 = (\mathcal O_c\T\mathcal O_c)\inv\mathcal O_c\T
        (z(0), \dot z(0), \dots)$. Conversely, if $\mathcal O_cx_0 = 0$ for some $x_0 \ne 0$, then $HF^kx_0 = 0$ for $k < n$, and by the Cayley--Hamilton theorem for all
        $k$, so $z(t) = \sum_k\frac{t^k}{k!}HF^kx_0 = 0$ for all $t$: $x_0$ produces the same output as the zero state, and the system is unobservable.

  \item Let the state be $(\theta, b_g)$ with the gyro reading $\omega_m$ as the known input: $\dot\theta = \omega_m - b_g$, $\dot b_g = 0$, and the inclinometer
        measures $z = \theta$. Then $F = \begin{bmatrix}0 & -1\\0 & 0\end{bmatrix}$, $H = [1\ \ 0]$, and $\mathcal O = \begin{bmatrix}1 & 0\\0 & -1\end{bmatrix}$ has full rank: the bias is
        observable, as the slope of the difference between the integrated gyro and the inclinometer. With a second gyroscope instead, the two
        readings are $\omega + b_1$ and $\omega + b_2$ of the same unknown rate $\omega$. Shifting $\omega \to \omega + c$, $b_1 \to b_1 - c$, $b_2 \to b_2 - c$ leaves both readings
        unchanged for every $c$, so only the difference $b_1 - b_2$ is observable: neither bias, nor the angle, can be recovered. Observing a bias
        needs an independent measurement of the quantity the sensor integrates, not another copy of the same sensor.
\end{solutions}

